% ---------------------------------------------------- methods, rebuilt from raw
\subsection{Adaptation against a checkpoint that is broken for a stated reason}
\label{sec:occ-methods}

The broken checkpoint in this task is broken for a named, traceable reason: it is the one
trained on the office where \texttt{rms} predicts occupancy, deployed where
\texttt{rms} predicts the opposite. That is the upgrade this section offers over reporting
a large adaptation gain against a weak model. The gain can be produced on demand, and
removed on demand, by two columns.

Table~\ref{tab:occmethods} runs the four methods under the Tier-1 constraints. The raw
files carry timestamps, so the target office is presented as a single pass in recording
order --- a requirement that could not even be stated on the anonymised extract, which has
no ordering to preserve --- with each batch predicted before the model updates on it, the
scaler fitted on the source split alone, and no target labels anywhere.

% source: src/exp35_occ_tta_raw.py -> results/exp35_occ_tta_raw.csv, 5 seeds
\begin{table*}[t]
\centering
\caption{Occupancy adaptation, MCC, Tier-1 single pass in recording order, $5$ seeds.
$\Delta$ is against source-only in the same cell. The final block is the prior extract of
Section~\ref{sec:occ-contrast}, retained for that contrast and not as a current result.}
\label{tab:occmethods}
\small
\begin{tabular}{llccccc}
\toprule
Features & Direction & Source-only & AdaBN & DIGA & PETAL & PeTTA \\
\midrule
\multicolumn{7}{l}{\emph{raw extracts}} \\
\addlinespace[2pt]
$9$ shared & source $\to$ target & $-0.3818$ & $\mathbf{+0.521}$ & $+0.075$ & $+0.519$ & $+0.509$ \\
$9$ shared & target $\to$ source & $+0.3579$ & $-0.119$ & $+0.011$ & $-0.159$ & $-0.168$ \\
$7$ transferable & source $\to$ target & $+0.4314$ & $+0.083$ & $+0.110$ & $+0.122$ & $+0.028$ \\
$7$ transferable & target $\to$ source & $+0.3465$ & $+0.046$ & $+0.019$ & $-0.009$ & $+0.005$ \\
\midrule
\multicolumn{7}{l}{\emph{prior extract, anonymised $9$-d --- see Section~\ref{sec:occ-contrast}}} \\
\addlinespace[2pt]
$9$-d & source $\to$ target & $+0.6628$ & $-0.022$ & $-0.020$ & $-0.025$ & $-0.080$ \\
$9$-d & target $\to$ source & $-0.0763$ & $+0.407$ & $+0.278$ & $+0.427$ & $+0.312$ \\
\bottomrule
\end{tabular}
\end{table*}

The pattern of Section~\ref{sec:res-mech} reproduces. Against the one checkpoint that is
actively anti-correlated with the labels, three methods gain roughly $52$ MCC points;
against every sound checkpoint they lose ground or barely move. What is new is the third
and fourth rows: restricted to the seven transferable features, no checkpoint is broken
and no method gains more than $12$ points. The adaptation gain and the feature defect
appear and disappear together, which is a stronger statement than the correlation between
gain and checkpoint quality that the activity task supports, because here the cause is
named.

\paragraph{A metric that rewards the wrong thing, for the third time.}
Section~\ref{sec:os-methods} showed PETAL recording the best open-set H-score while
classifying known activities half as well, because degrading a model uniformly raises the
component H was bottlenecked by. The same shape appears here in a different metric. Five
rows improve recall on the occupied class while losing MCC: AdaBN takes occupancy recall
from $0.355$ to $0.428$ in the target~$\to$~source cell while MCC falls from $+0.358$ to
$+0.239$, and TENT, PETAL and PeTTA do the same. Occupancy detection rate is a natural
headline for this task, and a paper reporting it alone would score all four as
improvements.

Three instances now, on three tasks, in three metrics: UCDR, H-score and occupancy recall.
In each the quantity of interest is a recall-like term that a method can raise by becoming
less decisive overall. The general form is worth stating: whenever a reported score is
dominated by one class's recall, it can be increased by a uniform degradation, and it
should not be published without the complementary term beside it.

\paragraph{Two ways to fail at $30\%$ positive, only one of which is usually checked.}
A binary task with roughly $30\%$ positives rewards predicting the majority class, and the
standard check is whether a method has collapsed to it. None here has: predicted-positive
rates stay between $17\%$ and $62\%$ across all twenty-four cells.

The failure that does occur is the opposite one. In the seven-feature
source~$\to$~target cell, the source-only model attains occupancy recall $0.967$ and
predicts occupied in $62\%$ of windows against a true rate of $28\%$ --- it has
over-committed to the \emph{minority} class. Its MCC of $0.431$ looks reasonable and its
accuracy of $0.640$ sits below the $0.720$ that predicting ``empty'' everywhere would
earn. Normalisation adaptation partly corrects it, which is why that row shows gains for
every method, and those gains are a calibration repair rather than a transfer improvement.
A collapse check that only looks for majority-class degeneracy passes this model without
comment, and an accuracy-only report would rank it below a constant predictor while its
MCC says otherwise. Both classes' recalls are needed to see what it is doing.
